% GENERATED FILE. Edit canonical lessons, then run npm run generate:books.
% canonical-source-hash: fnv1a64:aeb4e219e4d0cbd9
% canonical-lessons: LA-C71-maritus, LA-C71-uxor, LA-C71-liberi, LA-C71-infans, LA-C71-patruus

\chapter{Husband, Wife, Children}
\label{ch:la-husband-wife-children}
\hlchaptermodality{71}

\begin{chapteropening}
\textbf{By the end of this chapter:} \emph{I can name a husband, a wife, children, a baby and an uncle.}

Complete the last lesson of chapter 71: I can name a husband, a wife, children, a baby and an uncle.
\end{chapteropening}

\section[marītus, marītī]{marītus, marītī — a husband}

\label{lesson:LA-C71-maritus}



\begin{quote}
Before the new one: say the Latin for a grandmother, then the Latin for a family.
\end{quote}

\subsection*{You'll want to know: marītus, marītī}
\textbf{marītus, marītī} — \textquotedblleft{}a husband\textquotedblright{}. English \emph{marital}.

\begin{cousinweb}
\textbf{marītus} is probably inherited, from a Proto-Indo-European word for a young man of an age to marry, \emph{*meryo-}, \textquotedblleft{}young man\textquotedblright{} (EtymDB 2.1). English \emph{marry} comes from its family through French.
\end{cousinweb}

\begin{grammarlens}[title={What you've built}]
One more word for this chapter.
\end{grammarlens}

\subsection*{Your turn}
\begin{itemize}
\raggedright
  \item \emph{Say it:} \emph{marītus}
  \item \emph{Say it:} \emph{marītus}, once more
  \item \emph{Say it:} say \emph{marītus}
  \item \emph{Recall:} say the Latin for a grandmother, then the Latin for a family, then say \emph{marītus} again
\end{itemize}

\begin{tcolorbox}[breakable,colback=teal!4,colframe=teal!35!black,title={Before you move on}]
What does marītus, marītī mean? (\textquotedblleft{}A husband\textquotedblright{}.) Say it once more.
\end{tcolorbox}
\section[uxor, uxōris]{uxor, uxōris — a wife}

\label{lesson:LA-C71-uxor}



\begin{quote}
Before the new one: say the Latin for a family, then the Latin for a husband.
\end{quote}

\subsection*{You'll want to know: uxor, uxōris}
\textbf{uxor, uxōris} — \textquotedblleft{}a wife\textquotedblright{}.

\begin{cousinweb}
\textbf{uxor}, \textquotedblleft{}wife\textquotedblright{}, has no sure origin; Lewis \& Short call it doubtful, and no later proposal has won agreement. English \emph{uxorious}, \textquotedblleft{}doting on one's wife\textquotedblright{}, comes from it.
\end{cousinweb}

\begin{grammarlens}[title={What you've built}]
One more word for this chapter.
\end{grammarlens}

\subsection*{Your turn}
\begin{itemize}
\raggedright
  \item \emph{Say it:} \emph{uxor}
  \item \emph{Say it:} \emph{uxor}, once more
  \item \emph{Say it:} say \emph{uxor}
  \item \emph{Recall:} say the Latin for a family, then the Latin for a husband, then say \emph{uxor} again
\end{itemize}

\begin{tcolorbox}[breakable,colback=teal!4,colframe=teal!35!black,title={Before you move on}]
What does uxor, uxōris mean? (\textquotedblleft{}A wife\textquotedblright{}.) Say it once more.
\end{tcolorbox}
\section[līberī, līberōrum]{līberī, līberōrum — children}

\label{lesson:LA-C71-liberi}



\begin{quote}
Before the new one: say the Latin for a husband, then the Latin for a wife.
\end{quote}

\subsection*{You'll want to know: līberī, līberōrum}
\textbf{līberī, līberōrum} — \textquotedblleft{}children\textquotedblright{}. Only in the plural.

\begin{cousinweb}
\textbf{līberī}, \textquotedblleft{}children\textquotedblright{}, is the plural of \emph{līber}, \textquotedblleft{}free\textquotedblright{} (EtymDB 2.1): the free-born members of a household, as against its slaves. \emph{līber} is usually traced to Proto-Indo-European \emph{*h1lewdheros}, \textquotedblleft{}free\textquotedblright{} (EtymDB 2.1).
\end{cousinweb}

\begin{grammarlens}[title={What you've built}]
One more word for this chapter.
\end{grammarlens}

\subsection*{Your turn}
\begin{itemize}
\raggedright
  \item \emph{Say it:} \emph{līberī}
  \item \emph{Say it:} \emph{līberī}, once more
  \item \emph{Say it:} say \emph{līberī}
  \item \emph{Recall:} say the Latin for a husband, then the Latin for a wife, then say \emph{līberī} again
\end{itemize}

\begin{tcolorbox}[breakable,colback=teal!4,colframe=teal!35!black,title={Before you move on}]
What does līberī, līberōrum mean? (\textquotedblleft{}Children\textquotedblright{}.) Say it once more.
\end{tcolorbox}
\section[īnfāns, īnfantis]{īnfāns, īnfantis — a baby, an infant}

\label{lesson:LA-C71-infans}



\begin{quote}
Before the new one: say the Latin for a wife, then the Latin for children.
\end{quote}

\subsection*{You'll want to know: īnfāns, īnfantis}
\textbf{īnfāns, īnfantis} — \textquotedblleft{}a baby, an infant\textquotedblright{}. Literally \textquotedblleft{}not speaking\textquotedblright{}.

\begin{cousinweb}
\textbf{īnfāns} is built from \emph{in-}, \textquotedblleft{}not\textquotedblright{}, and \emph{for}, \textquotedblleft{}I speak\textquotedblright{} (Lewis \& Short): a baby is one who cannot yet speak. Its second part is \emph{fāns}, \textquotedblleft{}speaking\textquotedblright{}. English \emph{infant} comes from it.
\end{cousinweb}

\begin{grammarlens}[title={What you've built}]
One more word for this chapter.
\end{grammarlens}

\subsection*{Your turn}
\begin{itemize}
\raggedright
  \item \emph{Say it:} \emph{īnfāns}
  \item \emph{Say it:} \emph{īnfāns}, once more
  \item \emph{Say it:} say \emph{īnfāns}
  \item \emph{Recall:} say the Latin for a wife, then the Latin for children, then say \emph{īnfāns} again
\end{itemize}

\begin{tcolorbox}[breakable,colback=teal!4,colframe=teal!35!black,title={Before you move on}]
What does īnfāns, īnfantis mean? (\textquotedblleft{}A baby, an infant\textquotedblright{}.) Say it once more.
\end{tcolorbox}
\section[patruus, patruī]{patruus, patruī — an uncle (father's brother)}

\label{lesson:LA-C71-patruus}



\begin{quote}
Before the new one: say the Latin for children, then the Latin for a baby.
\end{quote}

\subsection*{You'll want to know: patruus, patruī}
\textbf{patruus, patruī} — \textquotedblleft{}an uncle (father's brother)\textquotedblright{}. Built on \textbf{pater}.

\begin{cousinweb}
\textbf{patruus} is built on \emph{pater}, \textquotedblleft{}father\textquotedblright{} (Lewis \& Short): the uncle on the father's side. \emph{pater} is usually traced to Proto-Indo-European \emph{*ph2ter}, \textquotedblleft{}father\textquotedblright{} (EtymDB 2.1).
\end{cousinweb}

\begin{grammarlens}[title={What you've built}]
One more word for this chapter.
\end{grammarlens}

\subsection*{Your turn}
\begin{itemize}
\raggedright
  \item \emph{Say it:} \emph{patruus}
  \item \emph{Say it:} \emph{patruus}, once more
  \item \emph{Say it:} say \emph{patruus}
  \item \emph{Recall:} say the Latin for children, then the Latin for a baby, then say \emph{patruus} again
\end{itemize}

\begin{tcolorbox}[breakable,colback=teal!4,colframe=teal!35!black,title={Before you move on}]
What does patruus, patruī mean? (\textquotedblleft{}An uncle (father's brother)\textquotedblright{}.) Say it once more.
\end{tcolorbox}
